\documentclass[11pt]{article}
\usepackage{icmlko}
\begin{document}

\title{압축은 애니 하나를 만들고 나머지 넷을 가리고 있었다}
\author{월드모델 실험실}
\date{노트 486 · 2026년 8월 3일}
\maketitle

\begin{abstract}
노트 485 가 다음 일을 이름까지 적어 뒀다 --- \emph{노트 $480 \cdot 481$
의 축 수 격차도 라벨 퍼짐 검사가 필요하다.} 거기서 나온 수들($+0.1028$ ·
도메인별 $+98\% \sim -67\%$)이 전부 압축에 오염됐을 수 있었다.

두 절반의 라벨 SD 는 \textbf{실제로 다르다}(차 평균 $0.129$ · 게임
$-0.262$ · 애니 $-0.172$). 그래서 맞춰서 다시 쟀다.

\begin{center}
\begin{tabular}{lrrr}
\hline
도메인 & 날 격차 & 압축이 만드는 격차 & 남은 것 \\
\hline
\textbf{애니} & $-0.1004$ & $\mathbf{-0.0897}$ & $\mathbf{-0.0108}$ \\
게임 & $+0.1180$ & $-0.0008$ & $+0.1188$ \\
세계애니 & $+0.1865$ & $-0.0090$ & $+0.1955$ \\
모바일 & $+0.2244$ & $-0.0521$ & $+0.2766$ \\
만화 & $+0.2640$ & $+0.0259$ & $+0.2382$ \\
\hline
\end{tabular}
\end{center}

\textbf{애니는 격차의 $89\%$ 가 압축이고 남는 것이 $-0.0108$ 뿐이다.}
나머지 넷에서는 압축이 격차를 \emph{만드는} 게 아니라 \textbf{가리고}
있었다 --- 빼면 격차가 커진다.

등록한 둘 중 ①(두 절반의 SD 가 다르다)은 통과, ②(맞추면 절반 이상
준다)는 \textbf{반증}이다.
\end{abstract}

\begin{figure}[t]
\centering
\includegraphics[width=\textwidth]{figs/f1.pdf}
\caption{\textbf{가} --- 도메인마다 셋. 파랑(압축)이 애니에서만 검정과
같은 방향이다. \textbf{나} --- 애니 이상의 종결.}
\end{figure}

\section{세 노트째 쫓던 애니가 닫힌다}

애니는 노트 $480 \cdot 481 \cdot 482$ 에서 계속 예외였다 --- \emph{축이
적은 행이 더 잘 맞는 유일한 도메인}. 노트 482 가 그 축(\texttt{venue\_
prominence})은 오히려 \emph{돕고} 있다는 것을 보였고, 노트 484 가
그 $119$행의 라벨이 좁다는 것을 찾았다($0.586$ 대 $0.721$).

\textbf{이 노트가 마지막 칸을 채운다.} 축 수로 가른 판에서도 같은 일이
일어나고, 퍼짐을 맞추면 $-0.1004$ 가 $-0.0108$ 로 \textbf{사라진다}.

\begin{quote}
\textbf{애니의 음수 격차는 라벨 압축이었다.} 축이 많은 행(=출처가 더
붙은 행)은 라벨이 좁은 무리이고, 좁은 라벨은 순위로 덜 갈린다. 축 자체는
그 행들을 $+0.051$ 돕고 있었다(노트 482).
\end{quote}

\section{나머지 넷에서는 반대로 걸린다}

모바일 $-0.0521$ · 세계애니 $-0.0090$ · 게임 $-0.0008$ --- 압축이 격차를
\textbf{줄이는} 방향으로 작용한다. 축이 많은 절반이 \emph{더 좁은} 라벨을
갖는 도메인들이라(게임 $1.809$ 대 $2.071$), 그 절반은 압축이라는 불리를
안고도 더 잘 맞힌 것이다. \textbf{빼고 보면 축의 값이 더 크다.}

\textbf{노트 481 의 수는 안 흔들린다} --- 오히려 그 도메인들에서 축의
몫은 적게 잡혀 있었다.

\section*{조항 --- 압축은 방향이 정해져 있지 않다}

노트 $484 \cdot 485$ 를 거치며 ``압축이 격차를 만든다''가 기본값처럼
됐는데, 이 노트가 그것을 막는다.

\begin{quote}
\textbf{조항} --- 라벨 퍼짐 검사는 \textbf{양방향}이다. 압축은 격차를
만들기도(\emph{애니}) 가리기도(\emph{모바일 · 게임}) 한다. \emph{어느
쪽인지는 두 무리의 SD 부호를 보고 정한다} --- 재기 전에 방향을 가정하지
않는다.
\end{quote}

그리고 노트 481 의 조항을 여기서 스스로 지킨다 --- \textbf{가중평균
($-40\%$)을 결론에 쓰지 않는다.} 부호가 갈렸다(애니만 양, 넷은 음).

\section*{남은 것}

\begin{itemize}
\item \textbf{모바일은 여전히 이상하다} --- 노트 481 에서 대조가 샌 게
아니었고 여기서도 압축이 아니다. 남은 설명은 노트 481 이 적은 \emph{대조
팔의 재학습}뿐이다.
\item \textbf{만화만 압축이 같은 방향}($+0.0259$)인데 크기가 작다. 만화는
노트 485 에서도 무리 안이 더 높은 유일한 도메인이었다.
\item \textbf{남은 $0.1439$(가중평균)} --- 압축을 뺀 뒤의 축 수 격차.
노트 480 의 대조(축을 줄여서 재기)와 이 대조(퍼짐 맞추기)는 서로 다른
것을 뺀다. 둘을 같이 빼면 얼마가 남는지 안 쟀다.
\end{itemize}

\end{document}
